\documentclass[11pt]{article}
\usepackage[a4paper, top=18mm, bottom=18mm, left=20mm, right=20mm]{geometry}
\usepackage[T2A]{fontenc}
\usepackage[utf8]{inputenc}
\usepackage[ukrainian]{babel}
\usepackage{graphicx}
\setlength{\parindent}{0pt}
\setlength{\parskip}{4pt}
\pagestyle{empty}
\begin{document}
%begin
{\large\textbf{Контрольна робота}}

\textbf{Варіант \No\,4}\hfill \textit{ПІБ:}\,\underline{\hspace{60mm}}
\vspace{4mm}

%beginitem
1. Що буде виведено за виконання фрагменту коду?

%code
print(9 != 4 >= 2 < False)
%code

%enditem

%beginitem
2. Що буде виведено за виконання фрагменту коду?
4**-1.0*True
%enditem

%beginitem
3. Відмітити все, що є коректним оператором С++:
( 1 ) n=1; n++

( 2 ) False=0

( 3 ) n=10; --n

( 4 ) x=float("inf")

%enditem

%end
\newpage
\end{document}

